\section{Introduction}


\begin{figure*}[ht]
  \centering
  \includegraphics[width=1.0\linewidth]{figures/overall.pdf}
  \caption{\textbf{Top right}: LeanDojo extracts proofs in Lean~\cite{de2015lean} into datasets for training machine learning models. It also enables the trained model to prove theorems by interacting with Lean's proof environment. \textbf{Top left}: The proof tree of a Lean theorem $\forall n \in \mathbb{N},~\texttt{gcd n n = n}$, where \hlblue{\texttt{gcd}} is the greatest common divisor (details in Sec.~\ref{sec:background}). When proving the theorem, we start from the original theorem as the initial state (the root) and repeatedly apply tactics (the edges) to decompose states into simpler sub-states, until all states are solved (the leaf nodes). Tactics may rely on premises such as \hlorange{\texttt{mod\_self}} and \hlgreen{\texttt{gcd\_zero\_left}} defined in a large math library. E.g., \hlorange{\texttt{mod\_self}} is an existing theorem $\forall n \in \mathbb{N},~\texttt{n \% n = 0}$ used in the proof to simplify the goal. \textbf{Bottom}: Our {\name} model (Sec.~\ref{sec:reprover}). Given a state, it retrieves premises from the math library, which are concatenated with the state and fed into an encoder-decoder Transformer~\cite{vaswani2017attention} to generate the next tactic.}
  \label{fig:overall}
\end{figure*}


Reasoning is a cornerstone of human intelligence and a fundamental goal of AI~\cite{newell1956logic}. One prominent task is automated theorem proving (ATP): automatically generating proofs for theorems expressed in formal logic. ATP is useful for formal mathematics, producing mathematical proofs that can be checked rigorously~\cite{buzzardfuture}. Furthermore, it underpins formal verification, which is essential for proving the correctness and safety of high-stakes applications~\cite{leroy2016compcert,ringer2019qed}.

ATP is challenging since the search space is prohibitively large. In many applications, it is impractical to generate proofs fully automatically. Therefore, interactive theorem proving (ITP) has emerged as an alternative paradigm. In ITP, proofs are constructed by human experts interacting with software tools called proof assistants, such as Coq~\cite{barras1997coq}, Isabelle~\cite{nipkow2002isabelle}, and Lean~\cite{de2015lean}.  
Machine learning can automate such interactive theorem proving, opening up a new avenue for theorem proving~\cite{yang2019learning}. The model can learn to interact with proof assistants, given data containing human-written proofs.

Formal theorem proving serves as an important challenge for machine learning. From a computer science perspective, formal proofs can be treated as programs~\cite{howard1980formulae}. But unlike conventional programs in C++ or Python, the correctness of proofs can be verified using proof assistants. Therefore, theorem proving may be considered a special form of code generation, with rigorous evaluation and no room for the model to hallucinate. This can be consequential to current large language models (LLMs), as they have demonstrated exceptional capability in code generation~\cite{chen2021evaluating} but have flaws in factuality and hallucination~\cite{ji2023survey}. In addition, augmenting LLMs with external tools, such as proof assistants, has shown promise in improving their various capabilities, including multi-step reasoning~\cite{schick2023toolformer}.

Current research on LLMs for theorem proving is facing many barriers. To our knowledge, none of the existing LLM-based provers are open-source~\cite{polu2020generative,jiang2021lisa,han2022proof,lample2022hypertree,jiang2022thor,polu2023formal,first2023baldur,wang2023dt}. They all use private pretraining data, and the compute requirements can reach thousands of GPU days~\cite{lample2022hypertree}. Furthermore, some rely on tailored infrastructure for distributed training and interaction with the proof assistant---both are not possible to fully reproduce without open-source code~\cite{polu2023formal,lample2022hypertree}. We change the status quo by introducing LeanDojo: open-source toolkits, models, and benchmarks that give researchers access to state-of-the-art LLM-based provers with modest computational costs.


\smallsec{Tools for Data Extraction and Interaction}
We focus on Lean, a proof assistant popular among mathematicians.\footnote{``Lean'' in our paper refers to Lean 3 by default. Lean 4 is not backward-compatible but is also supported by LeanDojo. Our Lean 4 results are in Appendix~\ref{appendix:lean4}.} Our framework LeanDojo provides two essential functions for learning-based theorem proving (Fig.~\ref{fig:overall}): extracting data and enabling models to interact with Lean programmatically. 


For data extraction, LeanDojo extracts training data not directly visible in the raw Lean code (Fig.~\ref{fig:lean_code}), e.g., proof trees consisting of intermediate states between proof steps (Fig.~\ref{fig:overall} \emph{Top left}). In addition, LeanDojo is the first tool to locate premises in Lean proofs, enabling training machine learning models for premise selection. For interaction, LeanDojo turns Lean into a gym-like interactive environment~\cite{brockman2016openai}. Using LeanDojo, the model can observe proof states, change the state by executing proof steps (referred to as ``tactics'' in proof assistants), and receive feedback from Lean. LeanDojo is the first tool capable of interacting with Lean reliably, reducing proof-checking errors in existing tools~\cite{polu2023formal} (correct proofs misjudged as incorrect) from 21.1\% to 1.4\%. 


\smallsec{Retrieval-Augmented LLMs for Theorem Proving}
LeanDojo addresses a key bottleneck in theorem proving: \emph{premise selection}~\cite{urban2004mptp,irving2016deepmath}. Existing LLM-based provers generate the next proof step (tactic), taking only the current state as input. However, proving theorems depends critically on the premises, such as lemmas and definitions, from a math library.

For example, Fig.~\ref{fig:overall} (\emph{Top left}) illustrates the proof of ``$\forall n \in \mathbb{N},~\texttt{gcd n n = n}$'', where \hlblue{gcd} stands for greatest common divisor. The proof starts from the original theorem as the initial state and repeatedly applies tactics to decompose states into simpler sub-states, until all states are solved. Tactics may rely on premises such as \hlorange{\texttt{mod\_self}} and \hlgreen{\texttt{gcd\_zero\_left}} defined in a large math library. E.g., \hlorange{\texttt{mod\_self}} is an existing theorem ``$\forall n \in \mathbb{N},~\texttt{n \% n = 0}$'' useful for simplifying the goal.

Incorporating all possible premises is too large to fit into LLMs' input, given the limited context window. Existing methods must learn to memorize the association between the proof state and the name \hlorange{\texttt{mod\_self}}. It works if the premise has been used in the training data to solve similar goals, but does not generalize to truly novel scenarios, e.g., theorems requiring lemmas unseen in training.    

One potential solution is to complement memorization with explicit premise selection. LeanDojo extracts premise data from Lean, including where they are defined and used. It enables us to tackle premise selection by augmenting LLMs with retrieval. We introduce \emph{\name} (\underline{Re}trieval-Augmented \underline{Prover}) (Fig.~\ref{fig:overall} \emph{Bottom}): Given the current state, it generates a tactic conditioning on a small number of premises retrieved from Lean's math library, \texttt{mathlib}~\cite{mathlib}. 
 
We need to limit retrieval to a small number of premises for it to be effective, and ideally, they should contain the ground truth premise. Our retriever builds upon Dense Passage Retriever (DPR)~\cite{karpukhin2020dense} but incorporates two algorithmic innovations: First, not all premises are accessible when proving a theorem (Sec.~\ref{sec:background}). LeanDojo can perform program analysis on Lean code to determine accessible premises. On our data, that reduces the average number of premises from 128K to 33K, significantly simplifying the retriever's task. Second, DPR needs negative examples in training and benefits from hard negatives, i.e., irrelevant premises that are hard to distinguish from ground truth ones. We propose in-file negatives: a simple mechanism to find hard negatives in premise selection, which samples negative premises defined in the same Lean source file as the ground truth premise. 


\smallsec{LeanDojo Benchmark} Using LeanDojo, we construct a benchmark containing {\numproofs} theorems/proofs extracted from \texttt{mathlib}. Our benchmark is one of the largest math-focused theorem-proving datasets. We find that the common practice of splitting theorems randomly into training/testing has led to an overestimated performance in the previous papers. LLMs can prove seemingly difficult theorems simply by memorizing the proofs of similar theorems during training. In LeanDojo Benchmark, we mitigate this issue by designing challenging data split requiring the model to generalize to theorems relying on novel premises that are never used in training.

We use LeanDojo Benchmark to train and evaluate {\name}. Training takes only five days on a single GPU. In evaluation, {\name} can prove 51.2\% theorems, outperforming a baseline that generates tactics directly without retrieval (47.6\%) and another baseline using GPT-4~\cite{openai2023gpt4} to generate tactics in a zero-shot manner (29.0\%). We also test {\name} on two existing datasets, MiniF2F~\cite{zheng2022minif2f} and ProofNet~\cite{azerbayev2023proofnet}. It can prove 26.5\% theorems in MiniF2F and 13.8\% in ProofNet, which is competitive with state-of-the-art methods without reinforcement learning~\cite{polu2023formal}, even though trained using far fewer resources. Moreover, it can prove 65 theorems that currently do not have proofs in Lean. Thus, our tool can also serve as an effective tool for augmenting existing math libraries in Lean.


\smallsec{Contributions}
In summary, we make four main contributions: First, we introduce tools for extracting data from and interacting with Lean. Second, we develop {\name}, the first retrieval-augmented language model for theorem proving. Third, we construct a challenging benchmark for learning-based theorem proving and use it to validate the effectiveness of {\name}. Finally, we facilitate open research on LLMs for theorem proving by releasing our data, model, and code. Our method does not rely on private datasets and can be trained on a single GPU within a week. We believe this will significantly lower the barriers to academic research in this area and establish the first accessible baselines for future work to build upon. Further, our method can be used to automatically generate new Lean proofs without requiring human effort.